\section{Erd\H{o}s Problem 42: reconstruction of the known general solution}

\subsection{1) Problem and source status}
Write $[N]=\{1,\ldots,N\}$. A set is Sidon if an equality
$a+b=c+d$ between its elements forces equality of the unordered pairs,
with repeated summands allowed. The intended assertion is that, for every
$M\geq1$ and all sufficiently large $N$, every nonempty Sidon set
$A\subset[N]$ admits a Sidon set $B\subset[N]$ such that
\[
 |B|=M,\qquad (A-A)\cap(B-B)=\{0\}.
\]
The nonempty qualification matters: if $A=\varnothing$, the intersection is
empty, and the appropriate formulation is containment in $\{0\}$.

The problem page checked on 23 September 2026 reports the general result as
solved, following a proof generated by GPT 5.5 Pro prompted by Harjas Sandhu
and subsequent discussion and formalization. The older local
\texttt{gpt\_pro\_5.4/42.tex} and \texttt{42\_v2.tex} cover finite
obstructions and the case $M=3$. This note gives a complete deduction of
the general result from the precisely stated compact Cayley theorem below.
That theorem is an external input; its substantial compactness proof is
not reproduced here. This is an exposition of existing work, with no
claim of a new solution or a newly checked Lean build.

\subsection{2) The external theorem and the plan}
For a function $h:\mathbb Z/p\mathbb Z\to\mathbb C$, use the normalized
Fourier transform
\[
 \widehat h(r)=\frac1p\sum_{x\bmod p}h(x)e^{-2\pi i rx/p}.
\]
We use the following theorem in the compact Cayley route to Problem 42.

\textbf{Compact Cayley theorem.}
For each integer $L\geq2$ and real $\eta>0$, there are
$\varepsilon>0$ and $p_0$ such that, for every prime $p>p_0$, the following
holds. If $T\subset\mathbb Z/p\mathbb Z$ satisfies
\[
 T=-T,\qquad 0\notin T,\qquad |T|\geq\eta p,
 \qquad \widehat{1_T}(r)\leq\varepsilon\quad(r\neq0),
\]
then there is a set $C$ of $L$ distinct residues with
\[
 c-c'\in T\quad\hbox{whenever }c,c'\in C,\ c\neq c'.
\]
The Fourier coefficients in this statement are real because $T=-T$.
The hypothesis is an upper bound, not a bound on their absolute values.
The exact theorem appears as
\texttt{Erdos42.CompactCayley.compact\_cayley\_clique} in the public
Lean source cited below. Its proof uses Fourier extraction, passage to a
connected compact group, and a counting argument. Reading that statement
and its downstream reduction does not constitute independently checking
all of the formal development.

We will put the allowed differences in $T$. The Sidon condition will give
the required Fourier upper bound. A translate of a sufficiently large
clique has many points in an interval of length $N$, and an elementary
greedy argument supplies a Sidon subset of those points.

\subsection{3) Greedy extraction, with all collision types included}
We first prove that every finite set $X\subset\mathbb Z$ with
$|X|\geq2M^3$ contains a Sidon subset of size $M$.

Suppose a Sidon subset $S\subset X$ of size $k<M$ has already been chosen.
An element $x\in X$ fails to extend it only if $x\in S$, or if a new
nontrivial sum equality has one of the forms
\[
 x+a=b+c\quad(a,b,c\in S),\qquad
 2x=a+b\quad(a,b\in S).
\]
Indeed, an equality with $x$ once on each side cancels to a trivial
equality, while $2x=x+a$ forces $x=a$. The number of forbidden integers
is therefore at most
\[
 k+k^3+\frac{k(k+1)}2<2M^3.
\]
Some point of $X$ remains available. Starting from the empty set and
repeating this argument produces $M$ points. The midpoint equation
$2x=a+b$ is necessary here because the Sidon convention permits equal
summands.

\subsection{4) The exact Fourier calculation}
Fix a nonempty Sidon set $A\subset[N]$, and write $a=|A|$.
Its positive differences are distinct, so
\[
 \frac{a(a-1)}2\leq N-1.
\]
By Bertrand's postulate choose a prime $p$ with $4N<p<8N$.
All sums of two members of $A$ lie in an interval of length less than $p$;
consequently $A$, reduced modulo $p$, is still Sidon. Put
\[
 D=A-A\pmod p,\qquad T=(\mathbb Z/p\mathbb Z)\setminus D.
\]
Then $T=-T$, $0\notin T$, and
\[
 |D|=a(a-1)+1\leq2N-1<p/2,
 \qquad |T|>p/2.
\]

Let convolution mean the unnormalized sum
$(f*g)(x)=\sum_y f(y)g(x-y)$ on the cyclic group. Every nonzero
directed difference has exactly one representation, whereas zero has
$a$ representations. Thus the following is an exact identity:
\[
 1_D=1_A*1_{-A}-(a-1)1_{\{0\}}.
\]
For $r\neq0$, Fourier transformation gives
\[
 \widehat{1_T}(r)
 =\frac{a-1}{p}-p\left|\widehat{1_A}(r)\right|^2
 \leq\frac{a-1}{p}
 \leq\frac{\sqrt{2N-2}}{4N}.
\]
The last bound tends to zero uniformly over every admissible $A$.
No equidistribution assumption about $A$ is used.

\subsection{5) Completing the general-$M$ argument}
Fix $M\geq1$, put $R=2M^3$, and apply the compact Cayley theorem with
$L=8R$ and $\eta=1/2$. Choose $N$ large enough that its prime
$4N<p<8N$ exceeds $p_0$ and that
$\sqrt{2N-2}/(4N)\leq\varepsilon$.
These conditions depend only on $M$. The preceding calculation supplies
a clique $C\subset\mathbb Z/p\mathbb Z$ of size $8R$ whose nonzero
differences all avoid $D$.

Average over translates $t\in\mathbb Z/p\mathbb Z$:
\[
 \frac1p\sum_t |C\cap(t+[N])|=\frac{|C|N}{p}>R.
\]
Hence some $t$ has at least $R$ clique points in $t+[N]$.
Define their integer lifts by
\[
 X=\{i\in[N]:t+i\pmod p\in C\}.
\]
Because $N<p$, these are distinct lifts, and $|X|\geq R$.
For distinct $x,y\in X$, the residue of $x-y$ belongs to $T$.
It follows in particular that the integer $x-y$ cannot belong to $A-A$.
The greedy extraction lemma supplies a Sidon subset $B\subset X$ of
size $M$. Its nonzero differences avoid $A-A$. Since both $A$ and $B$
are nonempty, zero belongs to both difference sets. This proves
\[
 (A-A)\cap(B-B)=\{0\}
\]
and completes the desired deduction for every fixed $M$.

\subsection{6) Checks and scope of the conclusion}
There are three distinctions worth making explicit. First, the large
clique itself need not be Sidon; the greedy extraction is a separate
step. Second, its vertices initially lie in a cyclic group; averaging
over interval translates and then lifting them is what puts $B$ in the
required integer interval. Third, a Sidon set can be maximal under
inclusion without having maximum cardinality. Replacing an arbitrary
$A$ by an inclusion-maximal set does not justify applying the Fourier
uniformity theorem for cardinality-extremal Sidon sets. The proof here
uses the uniform one-sided estimate instead.

The finite checks in \texttt{verification/solved\_42\_43\_checks.py}
check the exact modular difference identity and density estimate for
every nonempty Sidon subset of $[N]$, $1\leq N\leq12$.
They are consistency checks for the elementary reduction. They do not
prove the compact Cayley theorem or certify a threshold for arbitrary
$M$. The present argument is qualitative and makes no explicit bound
on $N_0(M)$.

\textbf{Outcome: SOLVED, as a reconstruction using a cited external
theorem.} The deduction after the compact Cayley theorem has no remaining
gap; that theorem is the stated mathematical dependency of this note.

\subsection{7) Sources and attribution}
\begin{itemize}
\item T. F. Bloom, problem page and discussion for Problem 42,
\texttt{https://www.erdosproblems.com/42} and
\texttt{https://www.erdosproblems.com/forum/discuss/42}, checked
23 September 2026. The discussion attributes the general solution to
GPT 5.5 Pro prompted by Harjas Sandhu, with subsequent exposition and
analysis by the contributors named there.
\item Public compact Cayley formal development,
\texttt{https://github.com/Shashi456/erdos-formalizations/blob/main/Erdos/P42/CompactCayley/Proof.lean},
in particular \texttt{compact\_cayley\_clique} and
\texttt{theorem\_1\_1\_from\_compact\_cayley}, source read
23 September 2026. The source declares only Lean/Mathlib core axioms
for these exposed theorems. No local type-checking is claimed here.
\end{itemize}

The estimate measures coverage of the intended nonempty-$A$ problem by
the reconstruction, with the explicitly cited theorem accepted. It does
not measure novelty or an independent formal verification.
\subsection{8) COMPLETION ESTIMATE}
\textbf{COMPLETION: 100\%}
